%% =============================================================================
%%  Etapa 7 — Síntese
%% =============================================================================
\chapter{Síntese}\label{cap:sintese}

\begin{objetivo}
Reunir o que foi observado numa ideia central, formular o princípio que o
texto ensina e, só então, perguntar como ele se aplica hoje. Os rascunhos
abaixo são propostas para revisar, não conclusões fechadas.
\end{objetivo}

\section{A ideia exegética}

Uma boa síntese tem \emph{sujeito} (de que o texto fala) e
\emph{complemento} (o que ele diz sobre isso).

\begin{nota}[Rascunho]
\textbf{Sujeito:} quem Deus declara favorecido e quem está em situação de
perigo, diante do reino anunciado por Jesus.\\
\textbf{Complemento:} os discípulos pobres, famintos, que choram e são
rejeitados por causa do Filho do Homem são declarados felizes, porque Deus
inverterá sua situação; os que já estão saciados, satisfeitos e aprovados
por todos recebem um lamento de advertência, porque já receberam por inteiro
o que buscavam.
\end{nota}

\begin{pergunta}
Reescreva a ideia exegética com suas palavras em no máximo duas frases.
Ela explica os quatro pares, a dobradiça \gr{πλήν} e a referência aos
profetas?
\end{pergunta}

\section{A ponte: do texto para hoje}

\textcite{duvall2012} propõem uma \enquote{jornada interpretativa} em cinco
passos: (1)~entender o texto na situação original; (2)~medir a distância entre
aquela situação e a nossa; (3)~formular o princípio teológico que atravessa
essa distância; (4)~conferir o princípio com o restante da Bíblia;
(5)~aplicá-lo à situação atual.

\begin{ebtabela}{@{}>{\raggedright}p{.3\linewidth} >{\raggedright\arraybackslash}X@{}}
\EBcab{Passo} & \EBcab{Neste texto}\\\midrule
Situação original & discípulos na Galileia, numa sociedade de honra e vergonha, com pobreza real\\
Diferenças & outro sistema econômico e social; leitores com outras experiências de pobreza e de riqueza\\
Continuidades & dependência de Deus; tentação da autossuficiência; preço de falar em nome de Deus\\
Princípio (rascunho) & Deus não mede a vida pela fartura nem pela aprovação presentes; quem depende dele pode esperar a inversão que ele promete, e quem se satisfaz com o que já tem corre o risco de nada mais receber\\
Conferir com a Bíblia & \rb{Lc 1:51–53; 12:13–21; 16:19–31; Tg 2:5; 5:1–6; 1Tm 6:17–19}\\
\end{ebtabela}

\section{Apropriação}

\begin{aplicacao}[Perguntas para a vida]
Onde procuro minha \enquote{consolação} de forma a já recebê-la por inteiro
agora? De quem busco a aprovação \enquote{de todos}? Que tipo de fome ou
de choro este texto me convida a não temer? Como as felicidades orientam o
modo como trato os pobres (\rb{14:12–14})?
\end{aplicacao}

\begin{aplicacao}[Minha aplicação]
\vspace{3cm}
\end{aplicacao}

\section{O que ficou em aberto}

\begin{itemize}
  \item Verificar as leituras \enquote{a conferir}: base textual da BJ e da
    NTLH, numeração da BP, texto e análise do siríaco.
  \item Decidir entre as leituras de \enquote{naquele dia} (v.~23).
  \item Comparar a síntese com pelo menos dois comentários de tradições
    diferentes.
\end{itemize}
